% RA-L final-accepted format per https://www.ieee-ras.org/publications/ra-l/
% ra-l-information-for-authors/ : letterpaper, 10pt, journal, twoside.
% (Initial PaperCept submissions use a separate ieeeconf.cls wrapper not
% distributed outside PaperCept/IEEE Author Center; this repo targets the
% final RA-L journal format directly since this is a working paper /
% eventual arXiv preprint, not a live PaperCept submission.)
\documentclass[letterpaper, 10pt, journal, twoside]{IEEEtran}
\usepackage{amsmath}
\usepackage{amssymb}
\usepackage{graphicx}
\usepackage{booktabs}
\usepackage{xcolor}
\usepackage{hyperref}
\usepackage{cite}
\usepackage{float}

\hypersetup{colorlinks=true, linkcolor=blue, urlcolor=blue, citecolor=blue}

\title{Proprioception as a Segmentation Prompt: Force-Guided Object Grounding for Learning from Demonstration}

% Author block withheld for double-anonymous review (IEEE RAS journal
% policy, effective 2026-01-01: names and affiliations must not appear
% anywhere in the manuscript, appendix, or supplementary material until
% after acceptance).

\begin{document}
\maketitle

\begin{abstract}
Learning-from-demonstration (LfD) pipelines that reason about kinematic
constraints or physical interaction depend on knowing which object in the
scene is relevant at each phase of a demonstration. This is typically
solved either by hand-labelling, by gripper-event heuristics that only
identify the grasped object, or by label-free visual object discovery. We
show empirically, across multiple real demonstrations, that label-free
visual discovery (DINOv2 self-attention combined with monocular or real
sensor depth) fails to isolate task-relevant objects in typical in-lab
manipulation scenes -- including when real, not estimated, depth is
available -- indicating that vision alone is insufficient. We propose
using the robot's own proprioception as the grounding signal instead:
gripper-state transitions localise the grasped object in time, and the
wrist force/torque wrench, via the intrinsic contact ray of Bicchi
\cite{bicchi1990}, geometrically constrains the location of the
contact-receiving object once camera calibration is available. Both cues
seed a frozen promptable video segmentation model (SAM~2
\cite{ravi2024sam2}) and propagate a role-tagged object mask across the
full demonstration. We report a working system validated on six
real-robot recordings spanning four distinct manipulation tasks and
three different sensor-availability profiles (full sensing, force-only,
gripper-only), including a first case of tracking three simultaneous
task-relevant objects through one demonstration. The geometric wrench-ray
grounding experiment -- the paper's core evaluation -- is fully
implemented and awaiting a direct hand-eye calibration pass (reading the
camera-to-sensor transform from CAD proved unreliable, Section~\ref{sec:calib});
we report the evaluation protocol and preliminary (dry-run, uncalibrated)
evidence that the signal is directionally correct.
\end{abstract}

\section{Introduction}

Learning from demonstration increasingly decomposes a demonstration into
phases governed by kinematic constraints \cite{overbeek2025frontiers} or
physical contact \cite{overbeek2025adrr}, both of which presuppose an
answer to a prior question: \emph{which object, in the camera frame, is
the one the constraint or contact concerns?} Most existing approaches
sidestep this entirely, defining constraints on end-effector pose alone,
or identify only the grasped object from gripper state and leave any
contact-receiving object unidentified.

A natural alternative is label-free visual object discovery: a frozen
foundation model proposes salient objects with no task-specific training.
We test this directly with DINOv2~\cite{oquab2023dinov2} self-attention
combined first with monocular depth (Depth-Anything-V2
\cite{yang2024depthanythingv2}) and then with \emph{real} sensor depth
from an RGB-D camera; in both cases the resulting saliency fails to
isolate the task-relevant object, across 16 (recording, event) pairs
spanning six recordings, at a mean intersection-over-union of only 0.160
against the true object (Section~\ref{sec:dado}) -- including on
non-specular objects where depth is measured cleanly, ruling out sensor
noise as the cause. This motivates a proprioceptive alternative.

\textbf{Contributions.}
\begin{itemize}
  \item \textit{Sensor-Robust Event Detection:} Grasp, release, and
        contact moments are localised from gripper and wrist
        force/torque signals, degrading gracefully to whichever signal
        remains when the other is absent, verified with zero
        behavioural change on full-sensing recordings against two
        gripper-only and one force-only recording.
  \item \textit{Geometric Contact Grounding:} The intrinsic contact ray
        of Bicchi~\cite{bicchi1990}, computed from the wrist wrench alone
        and projected into the camera frame via calibrated extrinsics,
        grounds the contact-receiving object without any learned
        proposal step -- implemented and validated in dry-run mode,
        pending a hand-eye calibration pass to report quantitative
        accuracy.
  \item \textit{Arbitrary-Object Generalisation:} The role-tagged
        tracking pipeline generalises from a fixed grasped/contact-
        receiver pair to an arbitrary object count, demonstrated with
        three independently seeded and bidirectionally propagated
        objects tracked across a full 1013-frame recording with no
        cross-object contamination.
  \item \textit{Label-Free Discovery Fails, Even With Real Depth:}
        DINOv2 attention $\times$ real sensor depth achieves a mean IoU
        of only 0.160 (range 0.044--0.352) against the true object across
        16 events, confirming that proprioceptive grounding is not
        replaceable by a stronger depth signal alone.
\end{itemize}

\section{Related Work}

\textbf{Gripper-event affordance learning.} Prior self-supervised
approaches exploit robot motion to bootstrap object segmentation without
manual labels -- e.g. Boerdijk \emph{et al.}~\cite{boerdijk2021gripper} use
optical-flow motion cues during a grasp to segment the held object from
the manipulator and background. These approaches identify \emph{that} an
object is grasped, continuously, from motion; they do not address which
object receives contact once grasped, nor do they use discrete
proprioceptive events (grasp/press) to trigger a single, targeted
segmentation prompt the way we do.

\textbf{Label-free / promptable segmentation.} Segment Anything
\cite{kirillov2023sam} and its video extension SAM~2
\cite{ravi2024sam2} provide high-quality, class-agnostic segmentation
given a point or box prompt, but supply no mechanism for choosing
\emph{where} to prompt. Depth-and-attention object discovery methods
(DADO-style: DINOv2 attention re-weighted by monocular depth) attempt
label-free prompt selection; we show this fails on our scenes even with
real depth substituted for the monocular estimate.

\textbf{Contact sensing.} Bicchi's intrinsic contact sensing
\cite{bicchi1990} recovers the point and direction of a single contact
from a wrist force/torque measurement alone, without vision. To our
knowledge this geometric signal has not previously been used as a prompt
for visual segmentation. Learned approaches instead predict force jointly
with a visual goal from text and image (ForceSight~\cite{collins2023forcesight}
predicts a text-conditioned visual-force goal for mobile manipulation);
our contact ray is not learned or predicted -- it is the measured wrench,
transformed geometrically, which requires no training data and applies
directly to any single-contact event.

\textbf{Dense tactile contact sensing.} Vision-based tactile sensors
provide dense contact geometry directly at the point of touch --
NormalFlow~\cite{huang2025normalflow} tracks full 6-DoF object pose
during contact from a GelSight-style sensor's surface normal field. This
requires a dedicated tactile-sensing fingertip and only observes the
object once the fingertip itself is in contact. Our contact ray uses a
single wrist-mounted force/torque sensor already present on the rig for
other reasons, and recovers a coarser 3D ray rather than a full 6-DoF
pose, but requires no additional sensing hardware and localises contact
anywhere on the manipulator, not only at an instrumented fingertip.

\textbf{Hand-eye calibration.} Recovering the fixed transform between a
rigidly-mounted eye-in-hand camera and the robot's proprioceptive frame
is a classical problem, solved in closed form by Tsai and
Lenz~\cite{tsai1989handeye} via the AX=XB formulation; our calibration
procedure uses their method directly. Reading this transform off CAD
drawings instead of measuring it directly proved unreliable at the
fidelity available to us (Section~\ref{sec:calib}), the practical case
for the direct calibration route standard in this literature.

\textbf{Constraint and interaction identification.} The host lab's
existing LfD layers identify kinematic constraints from pose
\cite{overbeek2025frontiers} and physical interactions from contact
\cite{overbeek2025adrr}; this work is intended as the visual grounding
layer beneath both.

\section{Method}
\label{sec:method}

\subsection{Event detection}

Given a merged, pose-synchronised recording, three event types are
detected: \emph{grasp} and \emph{release} from gripper-width threshold
crossings at the midpoint of the observed open/closed range, and
\emph{contact} (``press'') from the baseline-subtracted wrist
force-magnitude peak, where the baseline is the median force magnitude
over the first 10\% of the recording. When only one of the two sensors is
present, the detector degrades to that sensor's events alone with zero
change in behaviour when both are present, verified on real recordings
lacking each sensor in turn (Section~\ref{sec:multievent}). For
recordings with more than one grasp/release/contact cycle, peaks are
grouped into activity clusters by time-gap, and the contact search is
restricted to the window between the surrounding grasp and release
rather than reporting the single global-maximum peak, which we show
(Section~\ref{sec:multievent}) silently attributes the wrong phase's
contact on a multi-phase recording.

\subsection{Geometric grounding of the contact-receiving object}

At a contact event, the wrist wrench $(\mathbf{f}, \boldsymbol{\tau})$
measured in the sensor frame determines, under a single-contact
assumption, the contact ray's closest point to the sensor origin and its
direction:
\begin{equation}
  \mathbf{r}_0 = \frac{\mathbf{f} \times \boldsymbol{\tau}}{\lVert \mathbf{f} \rVert^2}, \qquad
  \hat{\mathbf{f}} = \frac{\mathbf{f}}{\lVert \mathbf{f} \rVert}
  \label{eq:bicchi}
\end{equation}
following Bicchi \cite{bicchi1990}. Transforming $\mathbf{r}_0$ and
$\hat{\mathbf{f}}$ through the sensor-to-base and base-to-camera
transforms and projecting through the calibrated camera intrinsics
yields a 2D ray in the image, which is used as a point (or line) prompt
for SAM~2 on the contact-receiving object -- geometrically, without any
learned object proposal step.

\subsection{Grasped-object grounding}

The grasped object is seeded either by a vision-only heuristic -- SAM
automatic mask generation scored by position, size, and border distance,
rejecting any candidate mask exceeding 40\% of frame area -- or, once
available, by projecting the end-effector position directly into the
image via the same calibration. The 40\% cutoff, tuned to one
recording's object scale, is shown in Section~\ref{sec:seedplacement} to
reject a correct but larger contact-receiving object on a different
scene, motivating the geometric alternative.

\subsection{Propagation and multi-object generalisation}

Each seeded object is propagated bidirectionally through the full
recording with SAM~2.1 Hiera-Large's video predictor, yielding a
per-frame mask. The sidecar representation generalises to an arbitrary
number of independently seeded objects rather than a fixed
grasped/contact-receiver pair, each with its own role label.

\section{System and Data}

Six real-robot recordings on a Franka Research 3 arm with a Franka Hand
gripper, Bota wrist force/torque sensor, and ZED RGB(-D) camera are used.
Table~\ref{tab:recordings} summarises sensor availability and scene
content; the recordings differ in which sensors were actively
publishing, exercising the graceful-degradation behaviour of the event
detector under real (not synthetic) conditions.

\begin{table}[t]
\centering
\caption{Recordings used in this work.}
\label{tab:recordings}
\begin{tabular}{lccp{3.3cm}}
\toprule
\textbf{Rec.} & \textbf{Grip.} & \textbf{F/T} & \textbf{Task} \\
\midrule
1 & \checkmark & \checkmark & pick / press / drop (2 objects) \\
2 & --         & \checkmark & multi-object contact (4 objects) \\
3 & \checkmark & --         & fastener insertion \\
4 & \checkmark & --         & fastener insertion \\
5 & \checkmark & \checkmark & multi-phase: latch / cube pick-press-drop / latch \\
6 & \checkmark & \checkmark & latch contact (short clip) \\
\bottomrule
\end{tabular}
\end{table}

Recordings 2--6 include registered depth in addition to RGB; for
Recordings 2--4 this was decoded from the raw sensor transport where the
standard export tooling could not process it, while Recordings 5--6 use
the host lab's own depth export directly. Recording 5 is notable for
containing three distinct interaction phases in one continuous capture
(a latch/hinge contact, then a cube pick-press-drop cycle, then a second
latch contact); a naive single-event detector that reports only the
global peak force in the recording silently attributes the wrong phase's
event to the grasp cycle it should describe (Section~\ref{sec:multievent}).

\section{Experiments}
\label{sec:experiments}

\subsection{Label-free visual discovery fails, even with real depth}
\label{sec:dado}

DINOv2 attention combined with monocular depth (Depth-Anything-V2) was
evaluated on Recording~1's event frames and fails to isolate either
task-relevant object -- attention is low and spatially uniform on these
in-lab scenes. To rule out monocular depth estimation error as the cause,
the same recipe (DINOv2 attention $\times$ real-depth$^{0.5}$, thresholded
at the 85th percentile) was re-run substituting real sensor depth, scored
quantitatively against the ground-truth object mask each event's
proprioceptive cue actually seeds (Section~\ref{sec:method}), across every
recording with both real depth and a propagated ground-truth mask
(Recordings 2--6): 16 (recording, event) pairs in total. Mean
intersection-over-union between the label-free proposal and the true
object is $0.160$ (range $0.044$--$0.352$, Table~\ref{tab:dado}), and the
attention signal itself stays diffuse and object-agnostic on every
event tested (Fig.~\ref{fig:dado}) -- real depth is measurably better
than a monocular estimate at recovering scene geometry, but does not
localise the task-relevant object reliably at any tested event, contact
or otherwise.

\begin{table}[t]
\centering
\caption{Label-free proposal (DINOv2 attn.\ $\times$ real depth) vs.\
ground truth, all recordings with real depth. IoU is intersection over
union between the thresholded proposal and the true object mask.}
\label{tab:dado}
\begin{tabular}{lccc}
\toprule
\textbf{Recording} & \textbf{n events} & \textbf{mean IoU} & \textbf{range} \\
\midrule
2 (4-object) & 5 & 0.182 & 0.112--0.268 \\
3 (fastener, grasp only) & 1 & 0.044 & -- \\
4 (fastener, grasp only) & 1 & 0.052 & -- \\
5 (cube phase) & 3 & 0.215 & 0.113--0.352 \\
5, earlier extraction\footnotemark & 3 & 0.178 & 0.092--0.286 \\
6 (latch) & 3 & 0.128 & 0.117--0.140 \\
\midrule
all & 16 & 0.160 & 0.044--0.352 \\
\bottomrule
\end{tabular}
\end{table}
\footnotetext{An earlier, separately-processed extraction of Recording
5's cube-manipulation phase, from before the host lab's native depth
export was available for this bag; included to show the result is not
an artefact of one specific extraction pipeline.}

\begin{figure}[t]
\centering
\includegraphics[width=\columnwidth]{figures/dado_vs_groundtruth_all_trials.png}
\caption{Label-free discovery (DINOv2 attention $\times$ real depth) vs.\
ground truth across representative events. DINOv2 attention (middle
column) stays diffuse and object-agnostic; the resulting thresholded
proposal (right column, red) rarely aligns with the true object outline
(green).}
\label{fig:dado}
\end{figure}

\subsection{Vision-only seed placement is unreliable across scenes}
\label{sec:seedplacement}

The vision-only heuristic used to seed the grasped/contact-receiver roles
scores SAM's candidate masks by position, size, and border distance, with
a 40\% frame-area cutoff tuned on Recording~1's object scale. On
Recording~2, the correct contact-receiving object exceeds this cutoff and
is rejected outright, causing the seed to fall on background clutter and
propagate an unstable mask (17--78\% of frame area) for the full
recording. Two
automated generalisations of the heuristic were attempted -- relative
candidate-size ranking, and a real-depth near-camera prior -- and both
failed: the first mis-selects a reflective background surface on
Recording~1 that it correctly avoided under the original heuristic; the
second is blind precisely on the specular surfaces where the original
failure occurs, since sensor depth dropout is worst there. This negative
result motivates the geometric grounding method of
Section~\ref{sec:calib} over further heuristic tuning.

\textbf{Point vs.\ box prompts.} A related, sharper failure mode was
found in how a seed is prompted to SAM~2, independent of how the seed
location is chosen. A single point prompt on a multi-coloured object
(a Rubik's cube) segments only the locally-contiguous coloured region
under it rather than the physical object; switching to a box prompt
around the whole object fixes this at the seed frame. Propagating both
prompt types across a full recording (Recording~5's cube phase,
975 frames) reveals a second, more severe failure: once the seed frame
is behind it, the box-prompted track degrades gracefully to
near-empty when the object leaves the frame (mean 3{,}243\,px across
469 object-absent frames), while the point-prompted track does not
recognise object absence at all and drifts onto unrelated background
--- mean 247{,}594\,px (48\% of the 960$\times$540 frame) across the
same 469 frames, peaking at 495{,}142\,px (95\% of frame) on a bare
pegboard panel (Table~\ref{tab:seedablation}). The two prompt types are
comparable while the object is genuinely present (66{,}075 vs.\
72{,}763\,px mean); the difference is entirely in how each degrades once
the object is gone.

\begin{table}[t]
\centering
\caption{Point vs.\ box seed prompt, propagated across Recording 5's
cube phase (975 frames, SAM~2).}
\label{tab:seedablation}
\begin{tabular}{lccc}
\toprule
\textbf{Prompt} & \textbf{Window} & \textbf{n frames} & \textbf{mean px (max)} \\
\midrule
box   & object present & 506 & 66{,}075 (80{,}462) \\
box   & object absent  & 469 & 3{,}243 (18{,}486) \\
point & object present & 506 & 72{,}763 (368{,}128) \\
point & object absent  & 469 & 247{,}594 (495{,}142) \\
\bottomrule
\end{tabular}
\end{table}

\subsection{Wrench-ray geometric grounding (protocol and CAD-sensitivity evidence)}
\label{sec:calib}

This is the paper's core planned evaluation. The method itself
(Eq.~\ref{eq:bicchi} through image projection) is fully implemented and
runs in a dry-run mode that reports base-frame geometry without
projecting whenever calibration is not marked trusted, so that no
placeholder calibration is ever silently used. At the contact event of
Recording~1, wrist wrench magnitude was 26.4\,N directed
$10.1^\circ$ from vertical in the base frame, consistent with a
predominantly top-down press onto the visible contact-receiving object;
the Bicchi contact point $\mathbf{r}_0$ lies 2.0\,cm from the wrist
origin, perpendicular to the force axis by construction, a physically
plausible contact offset for the demonstrated task.

\textbf{CAD-derived calibration is not yet trustworthy.} The host lab's
eye-in-hand camera bracket CAD left one translation ambiguous between two
candidate lens positions 63\,mm apart (the camera's own stereo baseline,
along its housing's local axis). Both candidates were evaluated by
running the full projection pipeline (live, not dry-run) on every
recording with a real force event and an already-propagated
contact-receiver mask (Recordings 1, 2, 6) and measuring the projected
ray's pixel distance to the mask centroid. Both candidates project
\emph{outside the visible frame} on all three recordings (551--760\,px
from the image centre against a 540\,px-tall image), independent of
which of the two is chosen -- indicating the CAD-derived
\emph{orientation} of the bracket, not only the ambiguous 63\,mm
translation, is unreliable. This is a stronger and more actionable
negative result than an unresolved translation would be: it rules out
CAD-reading as sufficient and motivates a direct hand-eye calibration
(checkerboard/ArUco correspondences across several arm poses,
\texttt{cv2.calibrateHandEye}) as the only remaining path to a trusted
extrinsic, which is in progress.

\textbf{Planned evaluation.} Once a trusted calibration is available, for
each contact event across all recordings with real depth (Recordings
2--6): project $\mathbf{r}_0$ and $\hat{\mathbf{f}}$ into the image,
intersect the ray with the scene using the recorded depth map, and
measure the 3D distance between the intersection point and the nearest
point on the true contact-receiving object's surface. We expect this to
report a single, low (sub-centimetre to few-centimetre) error figure per
event as the headline quantitative result.

\subsection{Generalisation to multiple simultaneous objects}

The fixed-pair role model was generalised to an arbitrary object count
and demonstrated on Recording~2 with three independently seeded and
bidirectionally propagated objects (Fig.~\ref{fig:multiobj}), each
tracked across the full recording (1013/1013 frames) with no
cross-object mask contamination and no unbounded mask growth.

\begin{figure}[t]
\centering
\includegraphics[width=\columnwidth]{figures/identify_depth_multi/overlays/f0786_1782835535683396106.png}
\caption{Three independently-seeded, independently-propagated objects
composed into one sidecar frame, generalising past a fixed
grasped/contact-receiver pair.}
\label{fig:multiobj}
\end{figure}

\subsection{Sensor-degradation robustness and multi-cycle detection}
\label{sec:multievent}

Event detection was verified to degrade gracefully and correctly across
all three sensor-availability profiles present in the recordings: full
sensing (Recordings 1, 5, 6), force-only (Recording~2, no gripper events
detected or attempted), and gripper-only (Recordings 3--4, no contact
event detected or attempted); Table~\ref{tab:eventdet} summarises all
six recordings. Each fallback path was regression-tested
against a recording with full sensing to confirm zero behavioural change
when both sensors are present.

Recording~5 additionally exposed a failure mode not present in the other
recordings: a single-event detector that reports only the global-maximum
force peak in the whole recording silently attributes the wrong phase's
contact to the grasp cycle it is meant to describe, since Recording~5
contains three distinct interaction phases (a latch contact, a cube
pick-press-drop cycle, a second latch contact) and the largest force peak
occurs in the first, unrelated phase. Restricting the press-event search
to the window between the detected grasp and release events (rather than
the full recording) recovers the correct, cube-related contact event.
This failure is silent -- it produces a plausible-looking but wrong
event with no error signal -- and would not have been caught without
manually inspecting the seeded frame.

Recording~6 exposes a second, independent silent failure. Its gripper
never actuates: the measured width spans 0.07999710--0.07999776\,m, a
total travel of $6.6\times10^{-7}$\,m that is entirely sensor noise.
Because the open/closed threshold is computed as the midpoint of the
observed width range, that midpoint falls \emph{inside} the noise band,
and the detector reports a grasp at $t{=}0.06$\,s and a release at
$t{=}7.66$\,s. The consequence propagates: since the contact search is
restricted to the interval between grasp and release, the spurious
window also displaces the contact event, which was reported at
$t{=}3.34$\,s against a true force peak of 11.15\,N at $t{=}5.08$\,s.
Requiring a minimum gripper travel of 1\,mm before believing any
transition removes both errors -- the guard sits roughly 24$\times$ below
the smallest genuine grasp in this dataset (0.024\,m) and three orders of
magnitude above the observed noise floor, so it separates the two cases
without tuning. With the guard active the detector correctly reports
contact only, at the true peak.

\begin{table}[t]
\centering
\caption{Event detection across all six recordings used in
this work. Recording~5 fails without the grasp-window restriction
(global-maximum force lands in the wrong interaction phase); Recording~6
fails without the minimum-travel guard (a gripper that never actuates
yields phantom grasp/release from sensor noise). Both are corrected in
Section~\ref{sec:multievent}.}
\label{tab:eventdet}
\begin{tabular}{lccl}
\toprule
\textbf{Recording} & \textbf{Grip.} & \textbf{F/T} & \textbf{Verdict} \\
\midrule
1              & \checkmark & \checkmark & pass \\
2              & --         & \checkmark & pass \\
3              & \checkmark & --         & pass \\
4              & \checkmark & --         & pass \\
5              & \checkmark & \checkmark & FAIL (single-event) / pass (windowed) \\
6              & \checkmark & \checkmark & FAIL (no span guard) / pass (guarded) \\
\bottomrule
\end{tabular}
\end{table}

\section{Limitations}

All recordings come from a single lab, a single robot arm, and a small
number of camera poses; generalisation across hardware setups is
untested. The sample size (six recordings, four distinct task types) is small for a
generalisation claim and is treated here as a system validation rather
than a statistically powered study. The paper's central quantitative
claim -- wrench-ray grounding accuracy -- is not yet reported, pending
a direct hand-eye calibration pass (Section~\ref{sec:calib}); this is
disclosed rather than approximated with placeholder values. Label-free
discovery was tested with one
attention-based method and one depth-conditioning scheme; other label-free
approaches may perform differently.

\section{Conclusion and Future Work}

Proprioceptive signals -- gripper state and wrist wrench -- ground
task-relevant object identification where label-free visual discovery
does not: DINOv2 attention $\times$ real depth achieves a mean IoU of
only 0.160 across 16 events, and a point-seeded SAM~2 track drifts onto
unrelated background at up to 95\% of frame area where a box-seeded
track degrades gracefully to near-empty instead. A geometric method for
grounding the contact-receiving object from wrist force/torque alone is
implemented and ready for evaluation once a direct hand-eye calibration
pass replaces the CAD-derived extrinsic ruled out in
Section~\ref{sec:calib}; this is the immediate next step. Beyond that,
integrating the resulting role-tagged object tracks into the host lab's
existing kinematic-constraint and physical-interaction identification
layers \cite{overbeek2025frontiers,overbeek2025adrr} is the intended path
to a complete perception-to-LfD pipeline.

\begin{thebibliography}{9}

\bibitem{bicchi1990}
A.~Bicchi, ``Intrinsic contact sensing for soft fingers,'' in
\emph{Proc. IEEE Int. Conf. Robotics and Automation (ICRA)}, 1990.

\bibitem{kirillov2023sam}
A.~Kirillov \emph{et al.}, ``Segment Anything,'' in \emph{Proc. IEEE/CVF
Int. Conf. Computer Vision (ICCV)}, 2023.

\bibitem{ravi2024sam2}
N.~Ravi \emph{et al.}, ``SAM 2: Segment Anything in Images and Videos,''
in \emph{Proc. Int. Conf. Learning Representations (ICLR)}, 2025.

\bibitem{oquab2023dinov2}
M.~Oquab \emph{et al.}, ``DINOv2: Learning Robust Visual Features without
Supervision,'' \emph{Trans. Mach. Learn. Res. (TMLR)}, 2024.

\bibitem{yang2024depthanythingv2}
L.~Yang, B.~Kang, Z.~Huang, Z.~Zhao, X.~Xu, J.~Feng, and H.~Zhao,
``Depth Anything V2,'' in \emph{Advances in Neural Information
Processing Systems (NeurIPS)}, vol.~37, 2024.

\bibitem{overbeek2025frontiers}
A.~H.~G.~Overbeek, D.~Dresscher, H.~van~der~Kooij, and M.~Vlutters,
``Versatile kinematics-based constraint identification applied to robot
task reproduction,'' \emph{Frontiers in Robotics and AI}, vol.~12, art.
1574110, 2025.

\bibitem{overbeek2025adrr}
A.~H.~G.~Overbeek, H.~van~der~Kooij, D.~Dresscher, and M.~Vlutters,
``Identifying physical interactions in contact-based robot manipulation
for learning from demonstration,'' \emph{Advanced Robotics Research},
e202500109, 2025.

\bibitem{boerdijk2021gripper}
W.~Boerdijk, M.~Sundermeyer, M.~Durner, and R.~Triebel, ``Self-supervised
object-in-gripper segmentation from robotic motions,'' in \emph{Proc.
Machine Learning Research (PMLR)}, 2021.

\bibitem{collins2023forcesight}
J.~A.~Collins, C.~Houff, Y.~L.~Tan, and C.~C.~Kemp, ``ForceSight:
Text-guided mobile manipulation with visual-force goals,'' \emph{arXiv
preprint arXiv:2309.12312}, 2023.

\bibitem{huang2025normalflow}
H.-J.~Huang, M.~Kaess, and W.~Yuan, ``NormalFlow: Fast, robust, and
accurate contact-based object 6DoF pose tracking with vision-based
tactile sensors,'' \emph{IEEE Robot. Autom. Lett.}, vol.~10, no.~1,
pp.~452--459, Jan. 2025.

\bibitem{tsai1989handeye}
R.~Y.~Tsai and R.~K.~Lenz, ``A new technique for fully autonomous and
efficient 3D robotics hand/eye calibration,'' \emph{IEEE Trans. Robot.
Autom.}, vol.~5, no.~3, pp.~345--358, Jun. 1989.

\end{thebibliography}

\end{document}
